\chapter{Réservation matière : renforcement MTO-2}\label{ch:mto2}

Chaque matière est décrite par une fiche (classe \jv{FicheMatiere}). Son champ
\jv{stockDisponible} représente le stock physique : la quantité réellement présente. MTO-2
ajoute à la fiche une table \jv{reservations}, qui associe à chaque commande la quantité
qu'elle a engagée. La quantité libre pour une nouvelle production s'obtient par soustraction :
\[
\text{disponibilité opérationnelle} = \text{stock physique} - \sum \text{réservations}.
\]
Une réservation ne modifie pas le stock physique ; elle modifie seulement la disponibilité.
\cref{fig:reservation} illustre ce point avec 200 unités de GPL : \jv{CMD\_1} en réserve
125, et la disponibilité tombe à 75 unités.

\begin{figure}[H]
  \centering
  \begin{tikzpicture}[font=\small]
    \draw[ink, fill=white] (0,0) rectangle (8,1.1);
    \draw[resc, fill=resc!30] (0,0) rectangle (5,1.1);
    \draw[mtoc, fill=mtoc!10, dashed] (5,0) rectangle (8,1.1);
    \node[font=\small\bfseries] at (2.5,0.55) {réservée à CMD\_1 : 125};
    \node[font=\small] at (6.5,0.55) {libre : 75};
    \node[above=2pt] at (4,1.1) {stock physique : 200 unités de GPL};
    \draw[{Stealth[length=2mm]}-{Stealth[length=2mm]}, ink] (0,-0.35) -- (8,-0.35);
    \node[below] at (4,-0.4) {\small une deuxième commande de 125 unités ne peut pas consommer la zone réservée};
  \end{tikzpicture}
  \caption{Stock physique de 200 unités, dont 125 engagées : la disponibilité est de 75 unités.}
  \label{fig:reservation}
\end{figure}


\section{Réservation atomique}

Une commande est réservée de manière atomique : soit toutes les matières de sa nomenclature
sont réservables, soit elle n'engage aucune matière. Une commande qui demande du GPL, des
bouteilles et des accessoires ne réserve pas le GPL seul si les bouteilles manquent. Cette
règle évite de bloquer de la matière sans pouvoir produire. La fonction
\jv{reserverMatierePourCommande} contrôle toutes les lignes avant d'écrire la moindre
réservation.

\section{Consommation contrôlée}

La consommation physique se fait unité par unité, au poste qui consomme la matière. La
fonction \jv{consommerMatieresPourUnite} identifie la commande à laquelle appartient l'unité,
puis décrémente à la fois le stock physique et la réservation de cette commande. Deux cas
refusent la consommation : la commande n'a pas de réservation suffisante, ou le stock
physique est insuffisant. Le système lève alors une erreur explicite, au lieu de ramener la
quantité à zéro.

Ce choix corrige un défaut de la version antérieure. Le modèle consommait la matière demandée
même lorsque le stock était insuffisant, en le ramenant à zéro. Des unités étaient fabriquées
sans matière disponible, sans signalement. Le \cref{ch:cas} rapporte cette observation.

\section{Invariants matière}

Quatre invariants sont contrôlés à chaque clôture de commande, et leur résultat est journalisé :
\begin{itemize}
  \item le stock physique n'est jamais négatif ;
  \item la somme des réservations ne dépasse pas le stock physique ;
  \item une production ne démarre pas sans matière engagée ;
  \item aucune réservation ne reste orpheline à l'état terminal de la commande.
\end{itemize}
Le journal \jv{[INVARIANT MATIERE] cloture CMD\_1 : OK} atteste que ces contrôles sont satisfaits.
